\begin{pastexam}[label={px:ntu-108-dsaa-08}]{臺大資工所 108 學年度 第 8 題}
（填空）

(8\%) Given a directed graph G where vertex u (in G) has at least one path to any other vertex
(in G), the following pseudocode is designed to output an acyclic graph. However, in Lines 8, 10,
11, and 14, each of A, B, C, and D should be 0, 1, or 2. Write down the correct values of A, B, C,
and D so that the pseudocode can successfully output an acyclic graph.
\begin{lstlisting}[numbers=left, numberstyle=\footnotesize, basicstyle=\ttfamily\small,
    keywordstyle=\bfseries, morekeywords={while,if,else,end}, frame=single,
    columns=fullflexible, keepspaces=true, xleftmargin=2.5em, framexleftmargin=2em]
Cycle-Removal(G,u)
  Initialize S as a stack
  Label all vertices in G as Type-0
  Label u as Type-1
  S.push(u)
  while S is not empty
    v = the top vertex of S // Comment: not S.pop() here
    if v has a Type-A neighbor w in G
      Remove (v,w) from G
    else if v a Type-B neighbor w in G
      Label w as Type-C
      S.push(w)
    else
      Label v as Type-D
      S.pop()
    end if
  end while
  output G
\end{lstlisting}
\end{pastexam}

\begin{pastsoln}
答案：A $= 1$，B $= 0$，C $= 1$，D $= 2$。

原卷第 10 行少了 has（「v a Type-B neighbor」），意思和第 8 行一樣，讀成「v has a Type-B
neighbor w」。 「neighbor」指 $v$ 的出邊指到的頂點。

這是用堆疊寫的 DFS（資結書第~21.4.6~節）。 三種標記對應 DFS 的三種顏色：Type-0 是還沒發現
（白），Type-1 是已經發現、還在堆疊裡（灰），Type-2 是鄰居都處理完、已經彈出（黑）。 第 4 行
把起點 $u$ 標成 Type-1 再推進堆疊，所以推進堆疊的頂點都標成 Type-1，C $= 1$。 第 10--12 行是
「往下走到一個還沒發現的鄰居」，B $= 0$。 第 14 行是「鄰居都處理完，結束 $v$」，D $= 2$。 第 8 行
要刪的是指向\textbf{灰色}頂點的邊：灰色頂點都在堆疊裡，是 $v$ 的祖先（或 $v$ 自己），這種邊就是
\textbf{後向邊}（資結書第~21.5~節），A $= 1$。 有向圖有環若且唯若 DFS 有後向邊
（資結書定理~21.10），把後向邊全部刪掉，剩下的就沒有環。

\textbf{為什麼輸出一定無環。} 令 $v$ 被標成 Type-2 的先後順序為它的完成順序。
\begin{enumerate}
\item 每個頂點都會完成。 $u$ 到每個頂點都有路徑。 假設某個頂點 $x$ 從頭到尾都是 Type-0，取 $u$
到 $x$ 路徑上第一個從沒被發現的頂點 $x'$，它在路徑上的前一個頂點 $y$ 被發現過。 $y$ 被推進堆疊，
最後一定會走到第 14 行（每個頂點最多推一次，每條邊最多刪一次，迴圈會結束）；那時 $y$ 沒有
Type-0 的鄰居。 但邊 $(y, x')$ 只有在 $x'$ 是 Type-1 時才會被刪，$x'$ 又從沒被發現，所以這條邊還在，
$x'$ 是 $y$ 的 Type-0 鄰居，矛盾。
\item 留下來的邊 $(v, w)$，$w$ 一定比 $v$ 先完成。 $v$ 走到第 14 行時，第 8、10 行的條件都不成立：
$v$ 沒有 Type-1 也沒有 Type-0 的鄰居，所以它剩下的每個鄰居 $w$ 都已經是 Type-2。 $v$ 彈出後不會
再推進堆疊（只推 Type-0 的頂點），它的出邊以後也不會再被刪，這時剩下的就是輸出裡 $v$ 的出邊。
\item 若輸出有環 $v_1 \to v_2 \to \cdots \to v_k \to v_1$，由第 2 點，$v_2$ 比 $v_1$ 先完成、$v_3$ 比
$v_2$ 先完成……$v_1$ 比 $v_k$ 先完成，繞一圈得到 $v_1$ 比自己先完成，矛盾。
\end{enumerate}

\textbf{小例子。} 邊 $u \to a$、$a \to b$、$b \to u$、$u \to b$。 $u$ 推入；$u$ 有 Type-0 的鄰居 $a$，推入；
$a$ 推入 $b$；$b$ 的鄰居 $u$ 是 Type-1，刪掉 $b \to u$；$b$ 沒有鄰居了，標 Type-2 彈出；$a$ 的鄰居
$b$ 是 Type-2，$a$ 完成；$u$ 的鄰居 $a$、$b$ 都是 Type-2，$u$ 完成。 只刪了後向邊 $b \to u$，
$u \to b$ 這條前向邊留著，輸出 $u \to a \to b$、$u \to b$ 沒有環。

\textbf{另一組也會無環的值。} 嚴格來說 D $= 1$（完成的頂點仍標 Type-1）時輸出也沒有環：那樣第 8 行
會把指向任何已發現頂點的邊都刪掉，只剩 DFS 樹的樹邊，當然無環；上面的例子裡 $u \to b$ 也會被刪。
但它多刪了前向邊和交叉邊，三種標記也失去「白、灰、黑」的意義，題目要的是 D $= 2$ 這組：只刪
後向邊，其他的邊都留著。
\end{pastsoln}
